
\section{Introduction}
\label{sec:introduction}

Large Language Model (LLM) agents increasingly rely on external tools to answer questions and perform real-world tasks \cite{yao2023react, shen2023hugginggpt}.
As tool ecosystems grow from tens to hundreds or thousands of available APIs, a fundamental question arises: \emph{how should an agent determine which tools to invoke for a given query?}

This problem---\emph{tool routing}---is typically formulated as direct tool selection: given a natural language query, the LLM selects one or more tools from the full catalog \cite{li2025toolscope, qin2024toolllm}.
The formulation is natural and works for small catalogs.
But it requires the model to solve a compound problem in a single step: identify which tools are relevant, infer a valid multi-step composition, and verify that the resulting sequence is executable.
As catalogs grow, this compound problem becomes increasingly brittle.
Context windows fill with tool descriptions, the combinatorial space of possible tool sequences explodes, and the model has no structural signal to distinguish valid compositions from invalid ones.
The result is missed tool chains, hallucinated tool invocations, and outputs whose compositional validity cannot be verified without execution.

We argue that these difficulties are not inherent to tool routing itself, but arise from formulating it at the wrong level of abstraction.
Direct tool selection conflates two fundamentally different reasoning tasks:

\begin{itemize}
    \item \textbf{Semantic reasoning:} What domain entities does the user refer to?
    \item \textbf{Compositional reasoning:} What sequence of typed transformations connects those entities?
\end{itemize}

These tasks have different structure and are best solved by different mechanisms.
Semantic reasoning---mapping natural language to domain concepts---is a strength of language models.
Compositional reasoning---ensuring that tool outputs are type-compatible with downstream inputs---is a structural constraint that can be solved exactly by graph algorithms.
Treating both as a single unconstrained prediction forces the LLM to reason simultaneously about meaning and execution, without leveraging the structure that tool registries already provide.

A key observation motivates the reformulation.
User intent is expressed over domain objects, not tools.
A query like ``Get the logs for pods in the production namespace'' refers to entity types (\texttt{Namespace}, \texttt{PodLogs}) without naming any tool.
The tools that connect these entities (\texttt{list\_pods}, \texttt{get\_pod\_logs}) are implementation details that follow from the entity types involved.
If the model can identify the correct entity types, the corresponding tool chain follows deterministically from the composition graph.

This suggests a different formulation: \emph{Can tool routing be decomposed into semantic reasoning over entity types and compositional reasoning over a typed graph?}

We propose \textbf{Typed Composition Routing (TCR)}, a reformulation that makes this decomposition explicit (Figure~\ref{fig:approach}).
TCR rests on a core assumption: \emph{the tool ecosystem exposes an underlying typed transformation graph}, where each tool consumes an entity of one type and produces an entity of another (e.g., \texttt{get\_pod\_logs}: \texttt{Pod} $\rightarrow$ \texttt{PodLogs}).
This assumption holds naturally for typed API ecosystems---cloud platforms, infrastructure automation, enterprise services, and any domain with structured request/response schemas---where the composition graph can often be derived automatically from existing specifications (e.g., OpenAPI).
Domains without typed interfaces fall outside this formulation.

Tools form a directed \emph{composition graph} through shared entity types.
Given a query, the LLM predicts the source and target entity types---semantic reasoning.
Graph search then generates the set of structurally valid tool chains connecting those types---compositional reasoning.
For most queries, this produces exactly one chain; when multiple valid chains exist (e.g., \texttt{list\_projects} vs.\ \texttt{create\_project} for the same entity pair), intent-level selection among guaranteed-valid candidates is a separable concern.
Where direct routing maps a query to a tool sequence ($q \rightarrow \{t_1, \ldots, t_k\}$), TCR factorizes this into type prediction ($q \rightarrow (e_{\mathrm{src}}, e_{\mathrm{tgt}})$) followed by candidate generation ($(e_{\mathrm{src}}, e_{\mathrm{tgt}}) \rightarrow \{[t_1, \ldots, t_l]\}$).

\begin{figure*}[t]
\centering
\begin{tikzpicture}[
    node distance=0.5cm,
    box/.style={
        rectangle,
        draw,
        rounded corners=2pt,
        minimum width=2.8cm,
        minimum height=0.6cm,
        font=\small,
        align=center
    },
    arrow/.style={-{Stealth[length=4pt]}, thick},
    annot/.style={font=\footnotesize\itshape, text=gray!70!black}
]

% Traditional approach
\node[font=\small\bfseries] (trad-title) {Direct Tool Selection};

\node[box, below=0.4cm of trad-title, fill=blue!8] (query1) {Query};

\node[
    box,
    below=0.5cm of query1,
    fill=orange!10,
    minimum width=3.2cm,
    minimum height=1.0cm
] (llm1) {LLM\\{\scriptsize Select tools from catalog}};

\node[box, below=0.5cm of llm1, fill=red!8] (tools1) {Selected tools};

\draw[arrow] (query1) -- (llm1);
\draw[arrow] (llm1) -- (tools1);

\node[annot, below=0.15cm of tools1]
{no validity guarantee};

% Divider
\draw[dashed, gray!50, line width=0.5pt]
(4.2,0.5) -- (4.2,-4.5);

% Typed Composition Routing
\node[
    font=\small\bfseries,
    right=5.5cm of trad-title
] (our-title) {Typed Composition Routing};

\node[box, below=0.4cm of our-title, fill=blue!8] (query2) {Query};

\node[
    box,
    below=0.5cm of query2,
    fill=green!10,
    minimum width=3.5cm
] (types)
{
Semantic reasoning (LLM)\\
{\scriptsize Predict entity types}
};

\node[
    box,
    below=0.5cm of types,
    fill=yellow!8,
    minimum width=3.5cm
] (graph)
{
Compositional reasoning (Graph)\\
{\scriptsize Generate valid candidates}
};

\node[
    box,
    below=0.5cm of graph,
    fill=green!15
] (tools2)
{
Valid candidates
};

\draw[arrow] (query2) -- (types);
\draw[arrow] (types) -- (graph);
\draw[arrow] (graph) -- (tools2);

\node[annot, below=0.15cm of tools2]
{validity guaranteed within graph};

\end{tikzpicture}

\caption{
\textbf{Left:} Direct tool selection requires the LLM to select and compose tools from the full catalog in a single unconstrained step.
\textbf{Right:} Typed Composition Routing decomposes the problem: the LLM performs semantic reasoning (entity type prediction), while graph search generates the set of structurally valid tool chains connecting the predicted types.
For most queries, this set contains exactly one chain.
}

\label{fig:approach}
\end{figure*}

This paper investigates whether this reformulation leads to more reliable and scalable tool routing.
Our contributions are:

\begin{enumerate}

\item A \textbf{reformulation of tool routing} that decomposes it into semantic reasoning (entity type prediction) and structural candidate generation (graph search over a typed composition graph), reducing the LLM's role from tool selection to entity type classification (Section~\ref{sec:method}).

\item An \textbf{analytical recall decomposition} that separates end-to-end routing performance into entity type prediction accuracy and graph reachability, enabling failures to be attributed to either component (Section~\ref{sec:recall-decomp}).

\item \textbf{Empirical validation} across four LLMs, five tool domains, and 140 benchmark queries, demonstrating consistent improvements over direct tool selection (average $\Delta$F1 = +0.32) with zero hallucinated tool invocations (Section~\ref{sec:results}).

\item Evidence that the formulation \textbf{scales to production tool ecosystems}: evaluation on a production MCP server with 1,060 API operations, where the typed composition graph (50 entity types, 1,060 edges) is derived automatically from OpenAPI specifications without manual curation (Section~\ref{sec:production-eval}).

\end{enumerate}

